Blood collection tubes are among the most widely used single-use medical plastics: a large hospital processes tens of thousands per day, yet virtually all are incinerated alongside hazardous clinical waste. Manual sorting at scale poses unacceptable infection risk to workers, motivating an automated circular economy pipeline in which tubes are \emph{segregated} by cap type, \emph{opened}, \emph{cleaned} and \emph{sterilised}, then returned to \emph{reuse} or single-polymer material recovery. The tubes in this study are unused, uncontaminated specimens captured in a controlled laboratory setting. Their partially transparent, specularly reflective surfaces make colour-based classification less reliable under varying clinical lighting---a central motivation for the RGB-D (colour plus depth) approach taken in this work.

The target deployment is a conveyor-belt sorter presenting one tube at a time to an overhead camera, with a pneumatic air-jet directing each tube to the correct recovery stream. The system targets sub-25\,ms per-frame latency; all experiments are performed on static images in a controlled laboratory environment, establishing baselines before dynamic validation at belt speed.

A dataset of 3\,000 paired RGB and depth images was captured with an Intel RealSense D415 depth camera and annotated across four cap-closure classes (Universal, Push-on, Screwcap, Other) using SAM~2-assisted polygon labelling on Roboflow. Five YOLO-family instance segmentation models were evaluated across RGB, RGBD, and depth-only input modalities. Adding a depth channel to a compact 2.8\,M-parameter model (YOLO11n-RGBD, mask mAP\textsubscript{50--95}\,=\,0.929) outperformed a network ten times larger trained on RGB alone (YOLOv8m-seg, 0.905); YOLO26n-seg delivered the best speed--accuracy trade-off at 23.1\,ms per frame. A depth-only model reached 0.804, confirming that tube geometry is discriminative and robust to the reflective and transparent surfaces that challenge colour-based classification. A live inference pipeline demonstrates classification at 30--45\,fps, ready for integration with the conveyor-belt sorting hardware.

